\section*{Capítulo \ref{ch:b2:alpine-orogen} --- Anatomía de una cordillera: los Alpes}

\begin{solution}{exo:b2:alpine-orogen:1}
(a) $e = 1 - 27.0/36.0 = 0.25$, es decir, un \qty{25}{\percent}; $S = 36.0 - 27.0 =
\qty{9.0}{km}$. (b) La capa debe haber conservado su longitud (sin estiramiento ni adelgazamiento a lo largo de ella),
el corte debe ser paralelo a la dirección de transporte, y las líneas de referencia deben marcar lugares
donde las capas no se deslizaron unas respecto de otras.
\end{solution}

\begin{solution}{exo:b2:alpine-orogen:2}
(a) Un \omterm{def:b2:alpine-orogen:klippe}{klippe}: el gneis es un resto aislado de un \omterm{def:b2:alpine-orogen:nappes}{manto de corrimiento} que descansa sobre su autóctono;
la erosión ha eliminado el resto de la lámina a su alrededor. (b) Una \omterm{def:b2:alpine-orogen:klippe}{ventana tectónica}: la erosión ha
atravesado el manto de gneis y ha dejado al descubierto las rocas más jóvenes situadas bajo su cabalgamiento basal, que
por tanto buza alejándose de la ventana por todos lados.
\end{solution}

\begin{solution}{exo:b2:alpine-orogen:3}
(a) Complejo ultramáfico (el manto superior); complejo gabroico (la corteza oceánica inferior);
complejo de \omterm{def:b2:magma-emplacement:bodies}{diques} laminados (la corteza media); complejo volcánico en almohadillas (la corteza superior).
(b) El complejo de \omterm{def:b2:magma-emplacement:bodies}{diques} laminados. (c) La mayor parte de la corteza de los océanos amplios se consume por subducción;
las \omterm{def:b2:alpine-orogen:ophiolite}{ofiolitas} que sobreviven en tierra se formaron sobre todo en pequeñas cuencas oceánicas o entre arcos
insulares, y las tetisianas donde un continente se adelgazó hasta que se abrió un pequeño océano.
\end{solution}

\begin{solution}{exo:b2:alpine-orogen:4}
(a) Mantos del Monte Rosa y del Gran San Bernardo: el lado europeo (el Brianzonés y el
margen pasivo europeo); Zermatt--Saas y zona del Combin: el océano Piamonte--Liguria; sistema de la Dent
Blanche: el margen de Apulia. (b) Apulia, que acaba arriba. (c) La placa superior no
es subducida, y las láminas arrancadas de la placa inferior se apilan bajo ella en el orden en que
llegó la placa inferior: primero el océano lejano, al final y más abajo el margen próximo. Un plegamiento posterior del apilamiento,
o cabalgamientos fuera de secuencia, pueden alterar el orden.
\end{solution}

\begin{solution}{exo:b2:alpine-orogen:5}
(a) $2\pi G\,\Delta\rho = 2\pi\times6.674\times10^{-11}\times500 =
\qty{2.097e-7}{s^{-2}}$; $h = 1.70\times10^{-3}/2.097\times10^{-7} = \qty{8.1}{km}$.
(b) $2.097\times10^{-7}\times25\,000 = \qty{5.24e-3}{m.s^{-2}}$, es decir, unos
\qty{-524}{mGal}. (c) La raíz no es una placa infinita: es un cuerpo de algunas decenas de kilómetros
de profundidad y solo aproximadamente tan ancho como la cordillera, de modo que atrae menos a un punto de la superficie que una
placa del mismo espesor extendida hasta el infinito; cuerpos densos como la zona de Ivrea
añaden anomalías positivas, y el contraste de densidad en profundidad puede ser menor que el supuesto.
\end{solution}

\begin{solution}{exo:b2:alpine-orogen:6}
(a) $164 - 44 = \qty{120}{Myr}$. (b) $z = P/(\bar\rho g)$:
$2.8\times10^9/(3000\times9.8) = \qty{95.2}{km}$; $2.8\times10^9/(3300\times9.8) =
\qty{86.6}{km}$. (c) «Más de \qty{90}{km}» queda dentro de este intervalo: la profundidad depende tanto
de la densidad media elegida para la sobrecarga como de la presión.
\end{solution}

\begin{solution}{exo:b2:alpine-orogen:7}
(a) Un flanco de longitud $L$ que buza con un ángulo $\theta$ abarca $L\cos\theta$ en horizontal; todos los flancos
tienen el mismo buzamiento, de modo que $l = l_0\cos\theta$ y $e = 1 - l/l_0 = 1 - \cos\theta$.
(b) $\cos\theta = 0.70$: $\theta = \qty{45.6}{\degree}$. (c) $1 - \cos\qty{40}{\degree} =
\qty{23.4}{\percent}$, como en la figura de los chevrones; $1 - \cos\qty{60}{\degree}
= \qty{50}{\percent}$.
\end{solution}

\begin{solution}{exo:b2:alpine-orogen:8}
(a) $l_0 = 118/2.0 = \qty{59.0}{km}$; $S = 59.0 - 42 = \qty{17.0}{km}$; $e = 17.0/59.0 =
\qty{28.8}{\percent}$. (b) $A_e = 118 - 42\times2.0 = \qty{34.0}{km^2}$; $d = 34.0/17.0 =
\qty{2.0}{km}$. (c) Sí: el \omterm{def:b2:alpine-orogen:decollement}{despegue} está en la base de la cobertera de \qty{2.0}{km}, bajo
su techo; la cobertera está despegada por su base.
\end{solution}

\begin{solution}{exo:b2:alpine-orogen:9}
(a) La raíz de Airy es $4.5h$: a \qty{46.5}{\degree} N, $4.5\times2.74 = \qty{12.3}{km}$, un
Moho a \qty{42.3}{km}; a \qty{45.55}{\degree} N, $4.5\times0.65 = \qty{2.9}{km}$, un Moho
a \qty{32.9}{km}. (b) Bajo el relieve alto el Moho está \qty{6.0}{km} menos profundo de lo que predice
Airy; bajo el borde de la llanura del Po está \qty{28.3}{km} más profundo. (c) La raíz no está
compensada localmente: está desplazada hacia el sur y es mayor de lo que exige el relieve.
La placa se flexiona en conjunto bajo las cargas que soporta (el soporte regional de un capítulo del volumen de tercer
año), y una de esas cargas, el cuerpo de Ivrea, no es topografía en absoluto.
\end{solution}

\begin{solution}{exo:b2:alpine-orogen:10}
(a) $5.1/11.3 = 0.45$: menos de la mitad. (b) $5.1 + 6.2 = \qty{11.3}{km}$. El techo del
Jurásico está por encima del \omterm{def:b2:alpine-orogen:decollement}{despegue} superior; los \qty{6.2}{km} transmitidos a lo largo de ese \omterm{def:b2:alpine-orogen:decollement}{despegue}
pasaron por debajo de los pliegues del corte y se consumieron en estructuras situadas más al norte, de modo que la
capa superior nunca se acortó esa cantidad dentro del corte. (c) Una sola capa guía registra
solo el acortamiento de las capas a las que pertenece; cada \omterm{def:b2:alpine-orogen:decollement}{despegue} por encima o por debajo de ella puede
llevar desplazamiento hacia dentro o hacia fuera del corte. Compénsense varias capas, y compruébese con las áreas.
\end{solution}

\begin{solution}{exo:b2:alpine-orogen:11}
(a) Para un desplazamiento de \qty{100}{km}, $\sqrt{200^2 + 100^2} = \qty{223.6}{km}$; para
\qty{50}{km}, $\sqrt{200^2 + 50^2} = \qty{206.2}{km}$; para \qty{150}{km},
$\sqrt{200^2 + 150^2} = \qty{250.0}{km}$. (b) $\arctan(100/200) = \qty{26.6}{\degree}$ respecto de
la dirección norte--sur, una línea de convergencia próxima a la N\qty{150}{\degree}E citada por Rosenberg y
colaboradores. (c) El desplazamiento es incierto en $\pm\qty{50}{km}$, lo que mueve la convergencia
total unos \qty{20}{km} en uno u otro sentido; la tercera cifra del resultado queda muy por debajo de
esa incertidumbre.
\end{solution}

\begin{solution}{exo:b2:alpine-orogen:12}
(a) De $210/12.9 = \qty{16.3}{Myr}$ a $210/8.5 = \qty{24.7}{Myr}$ (los \unit{km} divididos por
\unit{mm.yr^{-1}} dan millones de años). (b) Más de $4\times60 = \qty{240}{km}$; al
menos $240 - 60 = \qty{180}{km}$ de la cuenca están hoy bajo los cabalgamientos alpinos o han sido
erosionados. (c) La cuenca se forma delante de la cuña de cabalgamientos y migra con ella: su parte
interna está siendo cabalgada continuamente por los cabalgamientos que la cargan, y su relleno más antiguo queda atrapado
en los cabalgamientos, como la \omterm{def:b2:alpine-orogen:foreland}{Molasa} Subalpina.
\end{solution}

\begin{solution}{pb:b2:alpine-orogen:1}
\noindent\textbf{Parte I.}

\textbf{1.} $3.0/\cos\qty{35}{\degree} = \qty{3.66}{km}$; $2.5/\cos\qty{50}{\degree} =
\qty{3.89}{km}$; $4.0/\cos\qty{40}{\degree} = \qty{5.22}{km}$.

\textbf{2.} $8.00 + 2\times3.66 + 5.00 + 2\times3.89 + 6.0/\cos\qty{10}{\degree} +
2\times5.22 + 10.0/\cos\qty{5}{\degree} = 8.00 + 7.32 + 5.00 + 7.78 + 6.09 + 10.44 + 10.04 =
\qty{54.67}{km}$; con el tramo repetido, $l_0 = \qty{60.7}{km}$.

\textbf{3.} La capa situada bajo la lámina cabalgante es la misma capa, llevada bajo la lámina por el
cabalgamiento: antes del cabalgamiento estaba en continuidad con el resto. Omitirla haría que la capa
restituida fuera demasiado corta.

\textbf{4.} $e = 1 - 48.0/60.7 = 0.209$, es decir, un \qty{20.9}{\percent}; $S = 60.7 - 48.0 =
\qty{12.7}{km}$.

\textbf{5.} $60.7 - 60.2 = \qty{0.5}{km}$, $0.5/60.7 = \qty{0.8}{\percent}$: dentro de la
precisión de la medida; el corte es aceptable.

\textbf{6.} $60.7 - 54.0 = \qty{6.7}{km}$, un \qty{11}{\percent}: dos capas de una misma lámina
no pueden restituirse a longitudes tan distintas. El corte no sería admisible tal como está dibujado: una
estructura que falta o está mal dibujada, una línea de referencia mal situada, un movimiento fuera del plano
del corte, o un \omterm{def:b2:alpine-orogen:decollement}{despegue} entre las dos capas.

\noindent\textbf{Parte II.}

\textbf{7.} $146/2.4 = \qty{60.8}{km}$.

\textbf{8.} Difiere de \qty{60.7}{km} en \qty{0.1}{km}: la longitud y el área coinciden.

\textbf{9.} $A_e = 146 - 48.0\times2.4 = 146 - 115.2 = \qty{30.8}{km^2}$.

\textbf{10.} $d = 30.8/12.7 = \qty{2.43}{km}$, igual a los \qty{2.4}{km} de la cobertera con
la precisión de los datos: el \omterm{def:b2:alpine-orogen:decollement}{despegue} es la base de la cobertera.

\textbf{11.} $160/2.4 = \qty{66.7}{km}$, \qty{6.0}{km} más que la longitud de la capa. El engrosamiento
de las evaporitas de la base de la cobertera, que fluyen hacia los núcleos de los anticlinales, añade área
sin añadir longitud; la compensación de áreas con un espesor constante deja entonces de funcionar.

\textbf{12.} La cobertera se desplazó sobre un \omterm{def:b2:alpine-orogen:decollement}{despegue} situado en su base y el zócalo no
intervino. Bajo el cinturón, el techo del zócalo debería ser plano o buzar suave y
regularmente, sin estar roto por los cabalgamientos de encima.

\noindent\textbf{Parte III.}

\textbf{13.} $22 + 18 + 25 = \qty{65}{km}$; $S = 65 - 15 = \qty{50}{km}$.

\textbf{14.} $e = 50/65 = \qty{76.9}{\percent}$, frente al \qty{20.9}{\percent} del
cinturón de antepaís: los mantos están mucho más acortados.

\textbf{15.} $30 - 22 = \qty{8}{km}$ ($e = \qty{26.7}{\percent}$). El propio zócalo está
acortado: \omterm{def:b2:alpine-orogen:skin}{tectónica de piel gruesa}.

\textbf{16.} Cinturón de antepaís $12.7/(11 - 4) = \qty{1.81}{km/Myr}$, es decir,
\qty{1.81}{mm} al año; mantos $50/(30 - 20) = \qty{5.0}{mm}$ al año.

\textbf{17.} Los mantos (de 30 a \qty{20}{Ma}) se deformaron antes que el cinturón de antepaís (de 11 a
\qty{4}{Ma}): la deformación migró hacia fuera, hacia el antepaís. Es el orden de los
Alpes: el cabalgamiento de Glaris en el Oligoceno y el Mioceno, los macizos externos en el Mioceno,
el Jura en el Mioceno superior y el Plioceno.

\textbf{18.} $12.7 + 50 + 8 = \qty{70.7}{km}$.

\noindent\textbf{Parte IV.}

\textbf{19.} $\qty{0.5}{cm.yr^{-1}}\times\qty{40}{Myr} = 0.005\times4\times10^7 =
\qty{2.0e5}{m} = \qty{200}{km}$.

\textbf{20.} $200 + 18 = \qty{218}{km}$.

\textbf{21.} Solo se cuenta lo que se conserva o se infiere: partes del océano pueden haber
sido subducidas antes de \qty{85}{Ma} o más deprisa que la velocidad media, láminas de \omterm{def:b2:alpine-orogen:ophiolite}{ofiolita} pueden haber
sido erosionadas o llevadas fuera del corte, y una convergencia oblicua al corte introduce
fondo oceánico que el corte no ve.

\textbf{22.} Zona interna $60 - 20 = \qty{40}{km}$ ($e = \qty{66.7}{\percent}$); retrocuña
$70 - 48 = \qty{22}{km}$ ($e = \qty{31.4}{\percent}$).

\textbf{23.} $12.7 + 50 + 8 + 40 + 22 = \qty{132.7}{km}$; dentro del intervalo de las estimaciones
para los Alpes Centrales, de unos 125 a \qty{160}{km} de acortamiento colisional (antes de añadir
el desgarre).

\textbf{24.} Anchura restituida $153 + 132.7 = \qty{285.7}{km}$; corteza restituida
$285.7\times30 = \qty{8571}{km^2}$; corteza actual $153\times45 = \qty{6885}{km^2}$. Falta:
$8571 - 6885 = \qty{1686}{km^2}$, un \qty{19.7}{\percent} de la corteza restituida; ha sido
erosionada (y vertida en las \omterm{def:b2:alpine-orogen:foreland}{cuencas de antepaís}) o subducida con la corteza inferior en el
manto. Sin esas pérdidas la corteza tendría en promedio $8571/153 = \qty{56}{km}$ de espesor.
El acortamiento colisional total a través del corte es de
\textbf{\qty[text-series-to-math=true]{133}{km}}.
\end{solution}
